%% toc: P3 24 24.6
\chapter{Elementary Amalgamation and Expansion of the Language}
\label{sec:p03-w03-p03-006-amalgamation}

Two structures that look alike from a common set of elements can often be
placed together inside one larger structure, so that each sits there as an
elementary substructure and the common elements are identified.  Such an
\emph{elementary amalgam} turns a statement about two separate structures
into a statement about one.  The same construction works when one of the
structures carries extra symbols: a structure in a small signature can be
placed elementarily inside the reduct of a structure in a larger signature.
This chapter states these constructions, proves them in the annex, and
records what they need.

The input in every case is a pair of structures together with an
identification of some of their elements, and a hypothesis saying that the
identified elements satisfy the same formulas of a fixed kind on both
sides.  The conclusion is a structure into which both embed, with the
identification respected; the kind of formula in the hypothesis decides
whether the embeddings are elementary or only embeddings.  All proofs use
the diagrams of \cref{def:diagrams}, the diagram tests of
\cref{thm:diagram-tests} and first-order compactness,
\cref{thm:fol-compactness}.  The amalgamation theorem and its family and
disjoint forms follow the order of the lecture notes
\cite[Lecture~19, pp.~1--6]{Speirs2013Amalgamation}; the existential form,
the language-expansion form and the examples are given here with complete
proofs in Annex~\ref{sec:annex-p03-w03-p03-006-amalgamation}.

\section{Signatures, reducts and amalgams}
\label{sec:p03-w03-p03-006-setting}

\begin{hbconvention}{def:p03-w03-p03-006-reduct-conventions}{Reducts, expansions and named tuples}
Let \(\Lang\subseteq\Lang^+\) be first-order signatures, where every symbol
of \(\Lang\) has the same kind and arity in \(\Lang^+\).  For an
\(\Lang^+\)-structure \(\mathcal C\), its \emph{reduct}
\(\mathcal C\mathbin{\upharpoonright}\Lang\) has the same domain and keeps
only the interpretations of the symbols of \(\Lang\); \(\mathcal C\) is then
an \emph{expansion} of that reduct.  A map
\(g:\mathcal B\to\mathcal C\mathbin{\upharpoonright}\Lang\) from an
\(\Lang\)-structure is an \emph{\(\Lang\)-elementary embedding} when it is
an elementary embedding of \(\Lang\)-structures.

A \emph{named tuple} is a family \(\bar a=(a_i)_{i\in I}\) of elements
indexed by a set \(I\), which may be empty or infinite.  For an
\(\Lang\)-structure \(\mathcal B\) and an \(\Lang^+\)-structure
\(\mathcal C\) with named tuples \(\bar a\) in \(B\) and \(\bar c\) in
\(C\), both indexed by \(I\), we write
\[
 (\mathcal B,\bar a)\equiv_{\Lang}(\mathcal C,\bar c)
\]
when, for every \(\Lang\)-formula \(\varphi(x_{i_1},\ldots,x_{i_n})\) whose
free variables are indexed by \(i_1,\ldots,i_n\in I\),
\[
 \mathcal B\entails\varphi(a_{i_1},\ldots,a_{i_n})
 \quad\Longleftrightarrow\quad
 \mathcal C\entails\varphi(c_{i_1},\ldots,c_{i_n}).
\]
Only symbols of \(\Lang\) occur on either side.  When the two structures
have the same signature the subscript is dropped.  As in
\cref{thm:fol-compactness}, the Axiom of Choice is assumed whenever a
signature enlarged by constants for the elements of the structures is
uncountable.
\end{hbconvention}

An \(\Lang\)-sentence has the same truth value in \(\mathcal C\) and in
\(\mathcal C\mathbin{\upharpoonright}\Lang\), and the same holds for
\(\Lang\)-formulas at any assignment; satisfaction looks only at the
symbols that occur.  This fact is used throughout and is checked at the
start of Annex~\ref{sec:annex-p03-w03-p03-006-amalgamation}.

\begin{hbdefinition}{def:p03-w03-p03-006-elementary-amalgam}{Elementary amalgam}
Let \(f:\mathcal A\to\mathcal B\) and \(f':\mathcal A\to\mathcal C\) be
elementary embeddings of \(\Lang\)-structures.  An \emph{elementary
amalgam} of \(\mathcal B\) and \(\mathcal C\) over \(\mathcal A\) is an
\(\Lang\)-structure \(\mathcal D\) together with elementary embeddings
\(h:\mathcal B\to\mathcal D\) and \(k:\mathcal C\to\mathcal D\) such that
\(h\circ f=k\circ f'\).  The amalgam is \emph{disjoint} when, in addition,
\[
 h(B)\cap k(C)=h(f(A)).
\]
\end{hbdefinition}

A disjoint amalgam identifies only the elements that the two maps force
together.  Without disjointness, an amalgam may identify further elements of
\(B\) with elements of \(C\);
\cref{def:p03-w03-p03-006-dlo-example} exhibits both kinds over the same
data.

\begin{hbmethod}{def:p03-w03-p03-006-diagram-procedure}{Building an amalgam from diagrams}
Every amalgam below is obtained by the same five steps.
\begin{enumerate}
  \item \emph{Name the elements.}  Give each element of the structure that
        is to be extended a constant, as in \cref{def:parameter-expansion}.
        Give each element of the structure to be embedded a constant:
        the constant of its partner if it is identified with an element on
        the other side, and otherwise a new constant used nowhere else.
  \item \emph{Check the naming.}  Two elements that receive the same
        constant must be identified by the hypothesis.  Equality and
        inequality are atomic, so they are among the formulas the
        hypothesis transfers; this is the only side condition on the
        naming.
  \item \emph{Write the theory.}  Take the elementary diagram of the
        structure to be extended, together with the elementary diagram of
        the other structure when the embedding is to be elementary, or its
        atomic diagram when it is to be an embedding, both read through the
        naming.  Add any inequalities that the conclusion requires.
  \item \emph{Reduce finite satisfiability to one formula.}  A finite part
        of the theory mentions finitely many constants for unidentified
        elements.  Conjoin its sentences from the second diagram,
        replace those constants by variables, and quantify them
        existentially.  The result is a formula about identified elements
        of the kind the hypothesis transfers, and its witnesses on the
        extended side satisfy the finite part.
  \item \emph{Read off the maps.}  Compactness gives a model.  By
        \cref{thm:diagram-tests}, the interpretations of the constants are
        the required embeddings.  Renaming the elements of the model makes
        the embedding of the extended structure an inclusion.
\end{enumerate}
\end{hbmethod}

\section{The amalgamation theorems}
\label{sec:p03-w03-p03-006-theorems}

The basic theorem allows the larger structure to carry a larger signature.
The amalgamation theorem in one signature is its special case
\(\Lang^+=\Lang\).

\begin{hbtheorem}{thm:p03-w03-p03-006-expansion-amalgamation}{proved=annex:p03-w03-p03-006-expand}{Amalgamation into an expansion}
Let \(\Lang\subseteq\Lang^+\), let \(\mathcal B\) be an
\(\Lang\)-structure and let \(\mathcal C\) be an \(\Lang^+\)-structure.
Let \(\bar a\) in \(B\) and \(\bar c\) in \(C\) be named tuples indexed by
the same set \(I\), and suppose
\[
 (\mathcal B,\bar a)\equiv_{\Lang}(\mathcal C\mathbin{\upharpoonright}\Lang,\bar c).
\]
Then there are an \(\Lang^+\)-structure \(\mathcal D\) with
\(\mathcal C\preccurlyeq\mathcal D\) and an \(\Lang\)-elementary embedding
\(g:\mathcal B\to\mathcal D\mathbin{\upharpoonright}\Lang\) with
\(g(a_i)=c_i\) for every \(i\in I\).
\end{hbtheorem}

The hypothesis concerns only \(\Lang\)-formulas, while the conclusion
gives an elementary extension of \(\mathcal C\) in the full signature
\(\Lang^+\).  The extra symbols cost nothing because the elementary diagram
of \(\mathcal C\) is kept whole and only its \(\Lang\)-consequences are
compared with \(\mathcal B\).

\begin{hbcorollary}{thm:p03-w03-p03-006-elementary-amalgamation}{proved=annex:p03-w03-p03-006-elem}{Elementary amalgamation}
Both clauses concern structures \(\mathcal B\) and \(\mathcal C\) for one
signature \(\Lang\).
\begin{enumerate}
  \item If \(\bar a\) in \(B\) and \(\bar c\) in \(C\) are named tuples
        indexed by the same set with
        \((\mathcal B,\bar a)\equiv(\mathcal C,\bar c)\), then there are an
        elementary extension \(\mathcal D\succcurlyeq\mathcal C\) and an
        elementary embedding \(g:\mathcal B\to\mathcal D\) with
        \(g(\bar a)=\bar c\) coordinatewise.
  \item If \(f:\mathcal A\to\mathcal B\) and \(f':\mathcal A\to\mathcal C\)
        are elementary embeddings, then \(\mathcal B\) and \(\mathcal C\)
        have an elementary amalgam over \(\mathcal A\) in the sense of
        \cref{def:p03-w03-p03-006-elementary-amalgam}, in which the map
        on \(\mathcal C\) is an inclusion.
\end{enumerate}
\end{hbcorollary}

Clause~1 is the form stated in the lecture notes
\cite[Lecture~19, Theorem~1, p.~1]{Speirs2013Amalgamation}, where the
common elements form a set with a bijection between its copies.  The
hypothesis is about a set of elements, not a substructure: nothing requires
the named elements to be closed under the functions of either structure.

When the common part is an elementary substructure of both structures, the
amalgam can be chosen so that it identifies nothing else.

\begin{hbtheorem}{thm:p03-w03-p03-006-disjoint-amalgamation}{proved=annex:p03-w03-p03-006-disjoint}{Disjoint amalgamation over an elementary substructure}
Let \(\mathcal A\), \(\mathcal B\) and \(\mathcal C\) be
\(\Lang\)-structures with \(\mathcal A\preccurlyeq\mathcal B\),
\(\mathcal A\preccurlyeq\mathcal C\) and \(B\cap C=A\).  Then there are an
elementary extension \(\mathcal D\succcurlyeq\mathcal C\) and an elementary
embedding \(g:\mathcal B\to\mathcal D\) such that \(g(a)=a\) for every
\(a\in A\) and
\[
 g(B)\cap C=A.
\]
\end{hbtheorem}

The condition \(B\cap C=A\) is a matter of bookkeeping: elementary
embeddings \(f:\mathcal A\to\mathcal B\) and \(f':\mathcal A\to\mathcal C\)
can be replaced by inclusions into isomorphic copies of \(\mathcal B\) and
\(\mathcal C\) that meet exactly in \(A\), and the theorem then yields a
disjoint elementary amalgam over \(\mathcal A\).
The proof uses elementarity of \(\mathcal A\) in both structures to move a
finite covering statement from \(\mathcal B\) to \(\mathcal C\) through
\(\mathcal A\). The stronger algebraically closed-set version is
\cref{thm:p03-w03-p03-006-acl-amalgamation} below.
\Cref{def:p03-w03-p03-006-overlap-failure} shows why a restriction on
algebraic elements is needed.

\subsection{Disjoint amalgamation over an algebraically closed set}
\label{sec:p03-w03-p03-006-algebraically-closed}

The common elements need not form an elementary substructure. What matters
for disjointness is that they already contain the finite sets definable
over them. The following formulation is the algebraically closed-set
version of the theorem in
\cite[Lecture~19, Theorem~2, pp.~3--5]{Speirs2013Amalgamation}. Its proof
below supplies the simultaneous avoidance step in full.

\begin{hbdefinition}{def:p03-w03-p03-006-algebraic-closure}{Algebraic closure of a parameter set}
For \(X\subseteq M\), the set \(\operatorname{acl}_{\mathcal M}(X)\)
consists of the elements belonging to some finite subset of \(M\)
definable by a first-order formula with finitely many parameters from
\(X\). Thus \(b\in\operatorname{acl}_{\mathcal M}(X)\) means that
\(\mathcal M\entails\theta(b,\bar a)\), where \(\bar a\) is from
\(X\) and \(\{c\in M:\mathcal M\entails\theta(c,\bar a)\}\) is
finite. The set \(X\) is \emph{algebraically closed in \(\mathcal M\)}
when \(\operatorname{acl}_{\mathcal M}(X)=X\). Every member of \(X\)
is algebraic over \(X\), using its singleton definition. This is
model-theoretic algebraic closure in an arbitrary signature, not the
assertion that \(X\) is an algebraically closed field.
\end{hbdefinition}

\begin{hblemma}{thm:p03-w03-p03-006-finite-avoidance}{proved=annex:p03-w03-p03-006-finite-avoidance}{Simultaneous avoidance of a finite set}
Let \(m\geq1\) and let \(\bar y=(y_1,\ldots,y_m)\). Assume
\(\varphi(\bar y,\bar a)\) has its parameters \(\bar a\) in
\(X\subseteq M\). Suppose a tuple
\((b_1,\ldots,b_m)\) satisfies it in \(\mathcal M\), and every
\(b_i\) lies outside \(\operatorname{acl}_{\mathcal M}(X)\).
For every finite \(F\subseteq M\) there is a tuple satisfying
\(\varphi\) whose every coordinate lies outside \(F\).
\end{hblemma}

Infinite coordinate projections alone would not prove this lemma. In an
infinite set with a named element \(a\), the relation
\(y_1=a\lor y_2=a\) has two infinite projections, but every pair in it
meets \(\{a\}\). The lemma's hypothesis instead provides one tuple
all of whose coordinates are nonalgebraic over the parameters. Its proof
shows that if a definable family of tuples has any finite set meeting
every tuple, it has such a set consisting of algebraic elements.

\begin{hbtheorem}{thm:p03-w03-p03-006-acl-amalgamation}{proved=annex:p03-w03-p03-006-acl-amalgamation}{Disjoint elementary amalgamation over an algebraically closed set}
Let \(\mathcal B,\mathcal C\) be structures in one signature, with
\(B\cap C=X\), and let \(\bar x=(x)_{x\in X}\). Assume
\((\mathcal B,\bar x)\equiv(\mathcal C,\bar x)\) and
\(\operatorname{acl}_{\mathcal B}(X)=X\). There are an elementary
extension \(\mathcal D\succcurlyeq\mathcal C\) and an elementary
embedding \(g:\mathcal B\to\mathcal D\) fixing \(X\) pointwise such
that
\[
 g(B)\cap C=X.
\]
No closure of \(X\) under quantifiers or separate elementary-substructure
assumption is required. Algebraic closure in \(\mathcal C\) follows
from the hypotheses; it is not an additional assumption.
\end{hbtheorem}

An elementary substructure \(\mathcal A\preccurlyeq\mathcal B\)
contains every finite set definable over \(A\): if a formula has exactly
\(r\) solutions in \(\mathcal B\), the first-order statement asserting
exactly \(r\) distinct solutions has witnesses in \(\mathcal A\), and
these exhaust the solutions in \(\mathcal B\). Consequently the earlier
elementary-substructure theorem is a special case. The example with a
unique \(P\)-element in
\cref{def:p03-w03-p03-006-overlap-failure} fails the new hypothesis too:
that element is algebraic over the common set and lies outside it.


\begin{hbcorollary}{thm:p03-w03-p03-006-family-amalgamation}{proved=annex:p03-w03-p03-006-family}{Amalgamating a family}
Let \(X\) be a set and let \(\langle\mathcal B_j:j\in J\rangle\) be a
nonempty family of \(\Lang\)-structures with \(X\subseteq B_j\) for every
\(j\).  Write \(\bar x\) for the named tuple \((x)_{x\in X}\) and suppose
\((\mathcal B_j,\bar x)\equiv(\mathcal B_{j'},\bar x)\) for all
\(j,j'\in J\).  Then there are an \(\Lang\)-structure \(\mathcal D\) and
elementary embeddings \(g_j:\mathcal B_j\to\mathcal D\) such that
\(g_j(x)=g_{j'}(x)\) for all \(j,j'\in J\) and \(x\in X\).
\end{hbcorollary}

This is the form in which amalgamation is used to realize many types in one
elementary extension \cite[Lecture~19, Corollaries, pp.~5--6]{Speirs2013Amalgamation}:
each type over \(X\) is realized in some elementary extension of a fixed
structure, and the corollary puts all those extensions into one structure.

\section{Amalgamation that preserves only existential formulas}
\label{sec:p03-w03-p03-006-existential}

\begin{hbdefinition}{def:p03-w03-p03-006-existential-formula}{Existential and universal formulas}
An \(\Lang\)-formula is \emph{existential} if it has the form
\(\exists y_1\cdots\exists y_m\,\chi\) with \(\chi\) quantifier-free and
\(m\ge 0\), and \emph{universal} if it has the form
\(\forall y_1\cdots\forall y_m\,\chi\) with \(\chi\) quantifier-free.  For
an \(\Lang\)-theory \(T\), let \(T_\forall\) be the set of universal
\(\Lang\)-sentences \(\sigma\) with \(T\entails\sigma\).
\end{hbdefinition}

Weakening the hypothesis of
\cref{thm:p03-w03-p03-006-elementary-amalgamation} from all formulas to
existential formulas weakens the conclusion in exactly one place: the map on
one side becomes an embedding.

\begin{hbtheorem}{thm:p03-w03-p03-006-existential-amalgamation}{proved=annex:p03-w03-p03-006-exist}{Existential amalgamation}
Let \(\mathcal B\) and \(\mathcal C\) be \(\Lang\)-structures and let
\(\bar a\) in \(B\) and \(\bar c\) in \(C\) be named tuples indexed by the
same set \(I\).  The following are equivalent.
\begin{enumerate}
  \item For every existential \(\Lang\)-formula
        \(\varphi(x_{i_1},\ldots,x_{i_n})\) with \(i_1,\ldots,i_n\in I\), if
        \(\mathcal C\entails\varphi(c_{i_1},\ldots,c_{i_n})\) then
        \(\mathcal B\entails\varphi(a_{i_1},\ldots,a_{i_n})\).
  \item There are an elementary extension
        \(\mathcal D\succcurlyeq\mathcal B\) and an embedding
        \(g:\mathcal C\to\mathcal D\) with \(g(c_i)=a_i\) for every
        \(i\in I\).
\end{enumerate}
\end{hbtheorem}

The next result is the step of the \L o\'s--Tarski theorem,
\cref{thm:p03-los-tarski}, that produces a model around a structure.

\begin{hbcorollary}{thm:p03-w03-p03-006-universal-part}{proved=annex:p03-w03-p03-006-univ}{Models of the universal part}
Let \(T\) be an \(\Lang\)-theory and \(\mathcal B\) an \(\Lang\)-structure.
Then \(\mathcal B\entails T_\forall\) if and only if there is an embedding
of \(\mathcal B\) into some model of \(T\).
\end{hbcorollary}

\begin{hbexample}{def:p03-w03-p03-006-existential-failure}{An embedding that cannot be made elementary}
In the signature consisting of the order symbol \(<\) alone, let
\(\mathcal C=(\N,<)\) and \(\mathcal B=(\mathbb Z,<)\), and name the single element \(0\) on both sides.  The inclusion
\(\N\subseteq\mathbb Z\) is an embedding, so every existential formula true
of \(0\) in \(\mathcal C\) is true of \(0\) in \(\mathcal B\), and
\cref{thm:p03-w03-p03-006-existential-amalgamation} gives
\(\mathcal D\succcurlyeq\mathcal B\) and an embedding
\(g:\mathcal C\to\mathcal D\) with \(g(0)=0\); the inclusion itself is one
such pair.  No such \(g\) is elementary.  In \(\mathcal D\) the element
\(0\) has a predecessor, because \(\mathcal B\entails\exists y\,(y<0)\) and
\(\mathcal B\preccurlyeq\mathcal D\), while \(\mathcal C\entails\neg\exists
y\,(y<0)\).  The hypothesis transfers existential formulas from
\(\mathcal C\) to \(\mathcal B\) but not the universal formula
\(\forall y\,\neg(y<x)\).
\end{hbexample}

\section{Expanding the language}
\label{sec:p03-w03-p03-006-expansion}

For an \(\Lang^+\)-theory \(T\) and a subsignature \(\Lang\subseteq\Lang^+\),
write
\[
 T^{\Lang}=\{\sigma:\sigma\text{ is an }\Lang\text{-sentence and }T\entails\sigma\}.
\]
This is the part of \(T\) that can be stated without the extra symbols.

\begin{hbtheorem}{thm:p03-w03-p03-006-reduct-models}{proved=annex:p03-w03-p03-006-reduct}{Models of the restricted theory}
Let \(\Lang\subseteq\Lang^+\), let \(T\) be an \(\Lang^+\)-theory and let
\(\mathcal B\) be an \(\Lang\)-structure.  Then
\(\mathcal B\entails T^{\Lang}\) if and only if there are a model
\(\mathcal M\) of \(T\) and an \(\Lang\)-elementary embedding
\(\mathcal B\to\mathcal M\mathbin{\upharpoonright}\Lang\).  Moreover, if
\(\mathcal M_0\entails T\) and
\(\mathcal B\equiv\mathcal M_0\mathbin{\upharpoonright}\Lang\), then
\(\mathcal M\) can be chosen with \(\mathcal M_0\preccurlyeq\mathcal M\).
\end{hbtheorem}

So the models of \(T^{\Lang}\) are, up to isomorphism, exactly the
elementary substructures of reducts of models of \(T\).  Adding symbols and
axioms about them changes which \(\Lang\)-structures occur as reducts, but
every model of the \(\Lang\)-consequences sits elementarily inside one.

The lemma \cref{thm:p03-elementary-reduct-extension} is the special case of
\cref{thm:p03-w03-p03-006-expansion-amalgamation} in which the named tuple
lists the domain of the \(\Lang^+\)-structure and the \(\Lang\)-structure
is an elementary extension of its reduct; the renaming step of
\cref{def:p03-w03-p03-006-diagram-procedure} then turns the embedding into an
inclusion.  Alternating that lemma along two elementary chains gives the
Robinson joint consistency theorem,
\cref{thm:p03-robinson-joint-consistency}, and with it the interpolation
and definability results of \cref{thm:p03-craig-robinson-beth}.

\section{A worked example and a failure}
\label{sec:p03-w03-p03-006-examples}

\begin{hbexample}{def:p03-w03-p03-006-dlo-example}{Amalgamating dense orders over the rationals}
Work in the signature consisting of the order symbol \(<\) alone, and
let \(\mathcal A=(\mathbb Q,<)\).  Let
\(\mathcal B\) be \(\mathbb Q\cup\{\sqrt2\}\) with the order of the reals.
Let \(\mathcal C\) be \(\mathbb Q\cup\{r\}\) for a new point \(r\), ordered
so that \(q<r\) exactly when \(q<\sqrt2\), for \(q\in\mathbb Q\).  All three
are countable dense linear orders without endpoints, and
\(B\cap C=\mathbb Q\).

\emph{Elementarity.}  By \cref{thm:dlo-qe}, every formula is equivalent
modulo the theory of dense linear orders without endpoints to a
quantifier-free formula with the same free variables.  An inclusion of
models of that theory preserves and reflects quantifier-free formulas, so it
preserves and reflects every formula.  Hence
\(\mathcal A\preccurlyeq\mathcal B\) and \(\mathcal A\preccurlyeq\mathcal C\).

\emph{An amalgam that is not disjoint.}  Take \(\mathcal D=\mathcal B\),
\(h\) the identity, and \(k:\mathcal C\to\mathcal B\) fixing \(\mathbb Q\)
and sending \(r\) to \(\sqrt2\).  The map \(k\) is an order isomorphism,
hence elementary, and \(h\) and \(k\) agree on \(\mathbb Q\).  Here
\(h(B)\cap k(C)=B\), which is larger than \(\mathbb Q\).

\emph{A disjoint amalgam.}  Let \(\mathcal D\) be the order obtained from
\(\mathbb Q\) by inserting a copy \(E\) of \((\mathbb Q,<)\) into the cut
\(\{q:q<\sqrt2\}\mid\{q:q>\sqrt2\}\).  It is a dense linear order without
endpoints: between two points of \(E\) lies a point of \(E\); between a
rational \(q\) and a point of \(E\) lies a rational, because the cut has no
rational endpoint; and between two rationals lies a rational.  Choose
distinct \(e_0,e_1\in E\), send \(\sqrt2\) to \(e_0\) and \(r\) to
\(e_1\), and fix \(\mathbb Q\).  Both maps are inclusions of dense orders
without endpoints up to renaming, hence elementary, and the two images
meet exactly in \(\mathbb Q\).  This is the situation guaranteed by
\cref{thm:p03-w03-p03-006-disjoint-amalgamation}.

\emph{What cannot be changed.}  In every elementary amalgam over
\(\mathcal A\), the images of \(\sqrt2\) and \(r\) lie in the same cut of
\(\mathbb Q\), since the formulas \(q<x\) and \(x<q\) with rational
parameters are preserved.  An amalgam may keep them apart or identify them,
but it cannot separate them by a rational.
\end{hbexample}

\begin{hbwarning}{def:p03-w03-p03-006-overlap-failure}{Disjointness fails over a common set that is not elementary}
Let \(\Lang=\{P\}\) with \(P\) unary.  Let \(\mathcal B\) have domain
\(\N\) with \(P^{\mathcal B}=\{0\}\).  Choose a bijection \(\sigma\) from
\(\N\) onto a set \(C\) with \(\sigma(1)=1\) and \(\sigma(n)\notin\N\) for
\(n\ne 1\), and let \(\mathcal C\) have domain \(C\) with
\(P^{\mathcal C}=\{p\}\), where \(p=\sigma(0)\).  Then \(B\cap C=\{1\}\),
the one-element set \(\{1\}\) carries a substructure of both, and
\((\mathcal B,1)\equiv(\mathcal C,1)\) because \(\sigma\) is an isomorphism
fixing \(1\).  So
\cref{thm:p03-w03-p03-006-elementary-amalgamation} applies.  In any
elementary amalgam, \(\mathcal D\entails\exists x\,\forall y\,(P(y)\leftrightarrow y=x)\),
because \(\mathcal B\) satisfies this sentence; the images of \(0\) and of
\(p\) both satisfy \(P\), so they coincide.  The amalgam is therefore never
disjoint over \(\{1\}\).  The substructure on \(\{1\}\) is not elementary in
\(\mathcal B\): it fails \(\exists x\,P(x)\).  The element \(0\) is
definable without parameters, and an amalgam must identify definable
elements; the elementarity hypothesis of
\cref{thm:p03-w03-p03-006-disjoint-amalgamation} excludes this, since an
elementary substructure contains every element defined over it.
\end{hbwarning}

The theorems also need elementarity, or at least existential agreement, at
the level of the hypothesis.  Agreement on atomic formulas alone does not
permit an amalgam with elementary maps; the example following
\cref{thm:p03-elementary-reduct-extension} shows two elementary diagrams
that are jointly inconsistent over a common substructure.

\section{Sources}
\label{sec:p03-w03-p03-006-sources}

The elementary amalgamation theorem, its version over an algebraically
closed common set, the amalgamation of a family and its use for realizing
types follow Speirs's notes of Scanlon's Berkeley course
\cite[Lecture~19, pp.~1--6]{Speirs2013Amalgamation}; those notes take
Hodges' \emph{A Shorter Model Theory} as their course book.  The proofs in
the annex are written for this reference from the diagram method of
\cref{sec:diagrams}.
